\begin{ex}
Thước kẻ học sinh thông thường có vạch số bắt đầu từ số nào?
\shortans{0}
\loigiai{
Thước kẻ đo độ dài có vạch đo bắt đầu từ vạch số $0$.
\textbf{Đáp số:} $0$.
}
\end{ex}

\begin{ex}
Khi đo độ dài một cây bút chì bằng thước kẻ, ta đặt đầu bút chì thẳng với vạch số 0, đầu kia chỉ vạch số 13. Hỏi cây bút chì dài bao nhiêu xăng-ti-mét?
\shortans{13}
\loigiai{
Cây bút chì dài đúng $13$ cm.
\textbf{Đáp số:} $13$.
}
\end{ex}

\begin{ex}
Đặt một đoạn que tính từ vạch số 2 đến vạch số 8 trên thước kẻ. Hỏi que tính đó dài bao nhiêu xăng-ti-mét?
\shortans{6}
\loigiai{
Độ dài que tính là: $8 - 2 = 6$ (cm).
\textbf{Đáp số:} $6$.
}
\end{ex}

\begin{ex}
Đặt một cục tẩy từ vạch số 3 đến vạch số 7 trên thước kẻ. Hỏi cục tẩy dài bao nhiêu xăng-ti-mét?
\shortans{4}
\loigiai{
Độ dài cục tẩy là: $7 - 3 = 4$ (cm).
\textbf{Đáp số:} $4$.
}
\end{ex}

\begin{ex}
Bạn Lan ước lượng quyển sách dài 20 cm. Khi đo thực tế bằng thước kẻ thấy dài 22 cm. Hỏi độ dài thực tế dài hơn ước lượng bao nhiêu xăng-ti-mét?
\shortans{2}
\loigiai{
Độ dài thực tế hơn ước lượng là: $22 - 20 = 2$ (cm).
\textbf{Đáp số:} $2$.
}
\end{ex}

\begin{ex}
Bạn Hùng ước lượng gang tay dài 15 cm. Khi đo thực tế thấy dài 13 cm. Hỏi ước lượng dài hơn thực tế bao nhiêu xăng-ti-mét?
\shortans{2}
\loigiai{
Ước lượng dài hơn thực tế là: $15 - 13 = 2$ (cm).
\textbf{Đáp số:} $2$.
}
\end{ex}

\begin{ex}
Một cây thước dài 30 cm. Bạn Mai đo chiếc bàn học thấy dài bằng 3 lần cây thước đó. Hỏi chiếc bàn học dài bao nhiêu xăng-ti-mét?
\shortans{90}
\loigiai{
Chiều dài bàn học là: $30 + 30 + 30 = 90$ (cm).
\textbf{Đáp số:} $90$.
}
\end{ex}

\begin{ex}
Bề dày của cuốn sách Toán là 1 cm, cuốn từ điển dày 5 cm. Hỏi cuốn từ điển dày hơn cuốn sách Toán bao nhiêu xăng-ti-mét?
\shortans{4}
\loigiai{
Cuốn từ điển dày hơn là: $5 - 1 = 4$ (cm).
\textbf{Đáp số:} $4$.
}
\end{ex}

\begin{ex}
Một mẩu bút chì đặt từ vạch số 5 đến vạch số 12 trên cây thước kẻ. Hỏi mẩu bút chì đó dài bao nhiêu xăng-ti-mét?
\shortans{7}
\loigiai{
Chiều dài mẩu bút chì là: $12 - 5 = 7$ (cm).
\textbf{Đáp số:} $7$.
}
\end{ex}

\begin{ex}
Để đo chiều dài phòng học, bạn Nam dùng bước chân. Biết mỗi bước chân dài khoảng 40 cm. Nam bước 2 bước được quãng đường dài bao nhiêu xăng-ti-mét?
\shortans{80}
\loigiai{
Quãng đường 2 bước chân là: $40 + 40 = 80$ (cm).
\textbf{Đáp số:} $80$.
}
\end{ex}

\begin{ex}
Một que diêm dài khoảng 4 cm. Xếp 3 que diêm nối tiếp nhau thành một hàng thẳng. Hỏi hàng que diêm đó dài bao nhiêu xăng-ti-mét?
\shortans{12}
\loigiai{
Chiều dài hàng que diêm là: $4 + 4 + 4 = 12$ (cm).
\textbf{Đáp số:} $12$.
}
\end{ex}

\begin{ex}
Một chiếc lá bàng dài 16 cm, một chiếc lá bưởi dài 9 cm. Hỏi chiếc lá bàng dài hơn chiếc lá bưởi bao nhiêu xăng-ti-mét?
\shortans{7}
\loigiai{
Chiếc lá bàng dài hơn là: $16 - 9 = 7$ (cm).
\textbf{Đáp số:} $7$.
}
\end{ex}

\begin{ex}
Đoạn thẳng AB dài 15 cm. Trên đoạn thẳng AB lấy điểm C sao cho AC dài 9 cm. Hỏi đoạn thẳng CB dài bao nhiêu xăng-ti-mét?
\shortans{6}
\loigiai{
Độ dài đoạn CB là: $15 - 9 = 6$ (cm).
\textbf{Đáp số:} $6$.
}
\end{ex}

\begin{ex}
Một thanh sô-cô-la dài 18 cm gồm các ô vuông nhỏ dài 2 cm mỗi ô. Hỏi thanh sô-cô-la có bao nhiêu ô vuông?
\shortans{9}
\loigiai{
Vì $9 \times 2 = 18$ nên thanh sô-cô-la có $9$ ô vuông.
\textbf{Đáp số:} $9$.
}
\end{ex}

\begin{ex}
Cạnh của một khung ảnh hình vuông dài 15 cm. Hỏi tổng độ dài hai cạnh liền nhau của khung ảnh đó là bao nhiêu xăng-ti-mét?
\shortans{30}
\loigiai{
Tổng độ dài hai cạnh là: $15 + 15 = 30$ (cm).
\textbf{Đáp số:} $30$.
}
\end{ex}

\begin{ex}
Bạn An dùng sợi chỉ đo chu vi một cái cốc được 24 cm. Đo đường kính miệng cốc được 7 cm. Hiệu độ dài hai lần đo là bao nhiêu xăng-ti-mét?
\shortans{17}
\loigiai{
Hiệu độ dài là: $24 - 7 = 17$ (cm).
\textbf{Đáp số:} $17$.
}
\end{ex}

\begin{ex}
Một chiếc đinh sắt đặt từ vạch số 1 đến vạch số 6 trên thước kẻ. Chiếc đinh sắt dài bao nhiêu xăng-ti-mét?
\shortans{5}
\loigiai{
Chiều dài chiếc đinh sắt là: $6 - 1 = 5$ (cm).
\textbf{Đáp số:} $5$.
}
\end{ex}

\begin{ex}
Một chiếc kim khâu đặt từ vạch số 2 đến vạch số 5 trên thước kẻ. Chiếc kim khâu dài bao nhiêu xăng-ti-mét?
\shortans{3}
\loigiai{
Chiều dài chiếc kim khâu là: $5 - 2 = 3$ (cm).
\textbf{Đáp số:} $3$.
}
\end{ex}

\begin{ex}
Một sợi dây thừng dài 20 mét. Người ta cắt đi 8 mét. Hỏi sợi dây thừng còn lại dài bao nhiêu mét?
\shortans{12}
\loigiai{
Chiều dài còn lại là: $20 - 8 = 12$ (mét).
\textbf{Đáp số:} $12$.
}
\end{ex}

\begin{ex}
Bàn tay của bé xòe rộng một gang tay đo được 12 cm. Bé đo mép bàn được 5 gang tay như thế. Chiều dài mép bàn là bao nhiêu xăng-ti-mét?
\shortans{60}
\loigiai{
Chiều dài mép bàn là: $12 + 12 + 12 + 12 + 12 = 60$ (cm).
\textbf{Đáp số:} $60$.
}
\end{ex}

\begin{ex}
Một hộp sữa tươi hình hộp chữ nhật cao 12 cm, đáy rộng 4 cm. Chiều cao hơn chiều rộng đáy bao nhiêu xăng-ti-mét?
\shortans{8}
\loigiai{
Chiều cao hơn chiều rộng đáy là: $12 - 4 = 8$ (cm).
\textbf{Đáp số:} $8$.
}
\end{ex}

\begin{ex}
Bạn Bình có hai cây thước: thước thứ nhất dài 15 cm, thước thứ hai dài 20 cm. Nối hai cây thước lại được một đoạn dài bao nhiêu xăng-ti-mét?
\shortans{35}
\loigiai{
Tổng độ dài hai cây thước là: $15 + 20 = 35$ (cm).
\textbf{Đáp số:} $35$.
}
\end{ex}

\begin{ex}
Chiều dài chiếc bảng lớp học là 3 mét. Một mét bằng bao nhiêu xăng-ti-mét?
\shortans{100}
\loigiai{
Một mét bằng $100$ xăng-ti-mét.
\textbf{Đáp số:} $100$.
}
\end{ex}

\begin{ex}
Đoạn thẳng CD dài 14 cm. Điểm M là điểm chính giữa của đoạn thẳng CD. Hỏi đoạn thẳng CM dài bao nhiêu xăng-ti-mét?
\shortans{7}
\loigiai{
Vì M ở chính giữa và $7 + 7 = 14$ nên CM dài $7$ cm.
\textbf{Đáp số:} $7$.
}
\end{ex}

\begin{ex}
Ước lượng chiều dài chiếc xe ô tô đồ chơi là 10 cm, đo thực tế là 12 cm. Tính tổng độ dài ước lượng và thực tế.
\shortans{22}
\loigiai{
Tổng độ dài là: $10 + 12 = 22$ (cm).
\textbf{Đáp số:} $22$.
}
\end{ex}
